\documentclass[a4paper,oneside]{memoir}
% Standard preamble for literate programming documents
% Copy this file verbatim to doc/preamble.tex in new projects
%
% This preamble is shared across all literate programming projects.
% The main document (e.g., doc/main.tex) should % Standard preamble for literate programming documents
% Copy this file verbatim to doc/preamble.tex in new projects
%
% This preamble is shared across all literate programming projects.
% The main document (e.g., doc/main.tex) should % Standard preamble for literate programming documents
% Copy this file verbatim to doc/preamble.tex in new projects
%
% This preamble is shared across all literate programming projects.
% The main document (e.g., doc/main.tex) should \input{preamble} after
% \documentclass but before \begin{document}.

% Font setup.  Under pdfLaTeX these are REQUIRED: without [T1]{fontenc} a
% monospace [[code]] quote (or \texttt) cannot select a T1 typewriter shape and
% silently falls back to the proportional OT1 roman font -- i.e. monospace
% "disappears".  iftex guards them so this preamble stays safe when copied into
% a LuaLaTeX/XeLaTeX project (which already use Unicode fonts and need neither).
% Do NOT comment these out.  See latex-writing's references/unicode-and-fonts.md
% ("Symptom 2").
\usepackage{iftex}
\ifpdftex
  \usepackage[utf8]{inputenc}
  \usepackage[T1]{fontenc}
  \usepackage{lmodern}
\fi
\usepackage[swedish,british]{babel}
\usepackage{booktabs}
% The audit tables of `scholar sessions export -f table` are longtables
% with array column types (doc/methods/tables/).
\usepackage{longtable}
\usepackage{array}

\usepackage{noweb}
\noweboptions{shift,breakcode,longxref,longchunks}
% noweb names a chunk by its page and, only when that page holds two or
% more chunks, a letter (\@subpageref tests \2on<page>).  So whether a
% reference such as "papers 1865" reads "1865b" depends on where the
% pages break, and the references' width moves the breaks: the scholar
% document flipped between two layouts on every pass and never settled
% (#279).  With the letter always given, a reference keeps its length
% when its chunk is alone on the page, and the passes settle.
% The "used in" and "continued in" lists (built from \@subpagepair)
% keep noweb's form on purpose: always giving the letter there too made
% a list of more than 140 entries on page 158 change its length with
% the layout (its ranges such as 342-46 and its entries depend on the
% page breaks), and a build from clean still flipped between two
% layouts after 12 passes.
\makeatletter
\def\@subpageref#1#2#3\\{\nwthepagenum{#1}{#2}}
% A run of three or more pages in a "used in" list ("342--46") reaches
% \@nwhyperpagenum with an empty label: noweb.sty's \@pagesl and
% \@pagenumsl build such a range as {lo--hi}{}, and \@nwhyperpagenum
% links every entry to noweb.<label>, so pdfTeX warned "name{noweb.}
% has been referenced but does not exist" (#311).  A range now stays
% unlinked text; an entry with a label links as before, and the text,
% hence the layout, is the same.
\def\@nwhyperpagenum#1#2{%
  \ifx\relax#2\relax #1\else\nwhyperreference{#2}{#1}\fi}
\makeatother

% Verbose citation style: full reference on first use, short form on
% repeats.  This goes with memoir + didactic: didactic sets
% autocite=footnote and places footnotes in the margin, so \autocite
% yields the full reference as a margin note.  Without memoir+didactic
% verbose footnotes land at the page foot; use alphabetic/authoryear
% there instead.  Cite with \autocite in body text; inside \ltnote or any
% \marginpar use \parencite (a footnote inside a margin float is fatal:
% "Float(s) lost").
\usepackage[
  natbib,
  backend=biber,
  style=verbose,
  citestyle=verbose,
  singletitle=false,
  maxbibnames=99,
]{biblatex}
\addbibresource{bibliography.bib}

\usepackage[inline]{enumitem}
\setlist[enumerate]{label=(\arabic*)}

\usepackage[all]{foreign}
\renewcommand{\foreignfullfont}{}
\renewcommand{\foreignabbrfont}{}

%\usepackage{newclude}
\usepackage{import}

\usepackage[strict]{csquotes}
\usepackage[single]{acro}

\usepackage{subcaption}

\usepackage[noend]{algpseudocode}
\usepackage{xparse}

\let\email\texttt

%\usepackage[outputdir=ltxobj]{minted}
\usepackage{minted}
\setminted{autogobble}

\usepackage{pythontex}
\setpythontexoutputdir{.}
\setpythontexworkingdir{..}

\usepackage{fancyvrb}

\usepackage{amsmath}
\usepackage{amssymb}
\usepackage{mathtools}
\usepackage{amsthm}
\usepackage{thmtools}
\usepackage[unq]{unique}
\DeclareMathOperator{\powerset}{\mathcal{P}}

\usepackage[binary-units]{siunitx}

\usepackage{pgf}

\usepackage[colorlinks]{hyperref}
\usepackage[capitalize]{cleveref}
\crefalias{question}{item}
\crefname{question}{question}{questions}
\Crefname{question}{Question}{Questions}
\creflabelformat{question}{#2\textup{#1}#3}

% adapted from https://tex.stackexchange.com/a/30726/17418
\newcommand\chnote[1]{%
  \begingroup
  \renewcommand\thefootnote{}\footnote{#1}%
  \addtocounter{footnote}{-1}%
  \endgroup
}

\usepackage[marginparmargin=right]{didactic}

% Make noweb's hyperlinked identifiers safe in headings (chapter/section titles
% that contain a [[code]] quote).  With -index, such a quote weaves into a
% \nwlinkedidentq hyperlink; in a moving argument (the .toc, the running header,
% the PDF bookmark) hyperref then fails fatally (under newer TeX Live with
% errors like "Undefined control sequence \hyper@@link" / "\Hy@reserved@a", or a
% stack overflow from noweb's self-redefining \nwhyperreference).  Two fixes,
% both required, let the project keep [[...]] in headings instead of falling
% back to \texttt{...}:
%
% (1) Replace \nwhyperreference with noweb's robust hyperref form so it is
%     written literally into moving arguments instead of expanding (and
%     misfiring its terminating \global\let) there, yet still links in the body.
\DeclareRobustCommand\nwhyperreference[2]{\hyperlink{noweb.#1}{#2}}
%
% (2) Strip noweb's code-quote macros from PDF bookmark strings, so a [[code]]
%     quote in a heading yields clean bookmark text without hyperref warnings.
\pdfstringdefDisableCommands{%
  \def\nwlinkedidentq#1#2{#1}%
  \let\Tt\relax
  \let\nwendquote\relax
}

% Unicode character replacements for symbols not in Latin Modern fonts
% These appear in TUI code and documentation describing terminal output
\usepackage{newunicodechar}
\newunicodechar{✓}{\checkmark}
\newunicodechar{✗}{$\times$}
\newunicodechar{○}{\ensuremath{\circ}}
\newunicodechar{─}{\textemdash}
\newunicodechar{📝}{\texttt{[note]}}
\newunicodechar{🤖}{\texttt{[AI]}}
% LaTeX Error: Unicode character € (U+0080) not set up for use with LaTeX.
\newunicodechar{€}{\texttt{[0x80]}}
% LaTeX Error: Unicode character “ (U+0093) not set up for use with LaTeX.
\newunicodechar{“}{\texttt{[0x93]}}
 after
% \documentclass but before \begin{document}.

% Font setup.  Under pdfLaTeX these are REQUIRED: without [T1]{fontenc} a
% monospace [[code]] quote (or \texttt) cannot select a T1 typewriter shape and
% silently falls back to the proportional OT1 roman font -- i.e. monospace
% "disappears".  iftex guards them so this preamble stays safe when copied into
% a LuaLaTeX/XeLaTeX project (which already use Unicode fonts and need neither).
% Do NOT comment these out.  See latex-writing's references/unicode-and-fonts.md
% ("Symptom 2").
\usepackage{iftex}
\ifpdftex
  \usepackage[utf8]{inputenc}
  \usepackage[T1]{fontenc}
  \usepackage{lmodern}
\fi
\usepackage[swedish,british]{babel}
\usepackage{booktabs}
% The audit tables of `scholar sessions export -f table` are longtables
% with array column types (doc/methods/tables/).
\usepackage{longtable}
\usepackage{array}

\usepackage{noweb}
\noweboptions{shift,breakcode,longxref,longchunks}
% noweb names a chunk by its page and, only when that page holds two or
% more chunks, a letter (\@subpageref tests \2on<page>).  So whether a
% reference such as "papers 1865" reads "1865b" depends on where the
% pages break, and the references' width moves the breaks: the scholar
% document flipped between two layouts on every pass and never settled
% (#279).  With the letter always given, a reference keeps its length
% when its chunk is alone on the page, and the passes settle.
% The "used in" and "continued in" lists (built from \@subpagepair)
% keep noweb's form on purpose: always giving the letter there too made
% a list of more than 140 entries on page 158 change its length with
% the layout (its ranges such as 342-46 and its entries depend on the
% page breaks), and a build from clean still flipped between two
% layouts after 12 passes.
\makeatletter
\def\@subpageref#1#2#3\\{\nwthepagenum{#1}{#2}}
% A run of three or more pages in a "used in" list ("342--46") reaches
% \@nwhyperpagenum with an empty label: noweb.sty's \@pagesl and
% \@pagenumsl build such a range as {lo--hi}{}, and \@nwhyperpagenum
% links every entry to noweb.<label>, so pdfTeX warned "name{noweb.}
% has been referenced but does not exist" (#311).  A range now stays
% unlinked text; an entry with a label links as before, and the text,
% hence the layout, is the same.
\def\@nwhyperpagenum#1#2{%
  \ifx\relax#2\relax #1\else\nwhyperreference{#2}{#1}\fi}
\makeatother

% Verbose citation style: full reference on first use, short form on
% repeats.  This goes with memoir + didactic: didactic sets
% autocite=footnote and places footnotes in the margin, so \autocite
% yields the full reference as a margin note.  Without memoir+didactic
% verbose footnotes land at the page foot; use alphabetic/authoryear
% there instead.  Cite with \autocite in body text; inside \ltnote or any
% \marginpar use \parencite (a footnote inside a margin float is fatal:
% "Float(s) lost").
\usepackage[
  natbib,
  backend=biber,
  style=verbose,
  citestyle=verbose,
  singletitle=false,
  maxbibnames=99,
]{biblatex}
\addbibresource{bibliography.bib}

\usepackage[inline]{enumitem}
\setlist[enumerate]{label=(\arabic*)}

\usepackage[all]{foreign}
\renewcommand{\foreignfullfont}{}
\renewcommand{\foreignabbrfont}{}

%\usepackage{newclude}
\usepackage{import}

\usepackage[strict]{csquotes}
\usepackage[single]{acro}

\usepackage{subcaption}

\usepackage[noend]{algpseudocode}
\usepackage{xparse}

\let\email\texttt

%\usepackage[outputdir=ltxobj]{minted}
\usepackage{minted}
\setminted{autogobble}

\usepackage{pythontex}
\setpythontexoutputdir{.}
\setpythontexworkingdir{..}

\usepackage{fancyvrb}

\usepackage{amsmath}
\usepackage{amssymb}
\usepackage{mathtools}
\usepackage{amsthm}
\usepackage{thmtools}
\usepackage[unq]{unique}
\DeclareMathOperator{\powerset}{\mathcal{P}}

\usepackage[binary-units]{siunitx}

\usepackage{pgf}

\usepackage[colorlinks]{hyperref}
\usepackage[capitalize]{cleveref}
\crefalias{question}{item}
\crefname{question}{question}{questions}
\Crefname{question}{Question}{Questions}
\creflabelformat{question}{#2\textup{#1}#3}

% adapted from https://tex.stackexchange.com/a/30726/17418
\newcommand\chnote[1]{%
  \begingroup
  \renewcommand\thefootnote{}\footnote{#1}%
  \addtocounter{footnote}{-1}%
  \endgroup
}

\usepackage[marginparmargin=right]{didactic}

% Make noweb's hyperlinked identifiers safe in headings (chapter/section titles
% that contain a [[code]] quote).  With -index, such a quote weaves into a
% \nwlinkedidentq hyperlink; in a moving argument (the .toc, the running header,
% the PDF bookmark) hyperref then fails fatally (under newer TeX Live with
% errors like "Undefined control sequence \hyper@@link" / "\Hy@reserved@a", or a
% stack overflow from noweb's self-redefining \nwhyperreference).  Two fixes,
% both required, let the project keep [[...]] in headings instead of falling
% back to \texttt{...}:
%
% (1) Replace \nwhyperreference with noweb's robust hyperref form so it is
%     written literally into moving arguments instead of expanding (and
%     misfiring its terminating \global\let) there, yet still links in the body.
\DeclareRobustCommand\nwhyperreference[2]{\hyperlink{noweb.#1}{#2}}
%
% (2) Strip noweb's code-quote macros from PDF bookmark strings, so a [[code]]
%     quote in a heading yields clean bookmark text without hyperref warnings.
\pdfstringdefDisableCommands{%
  \def\nwlinkedidentq#1#2{#1}%
  \let\Tt\relax
  \let\nwendquote\relax
}

% Unicode character replacements for symbols not in Latin Modern fonts
% These appear in TUI code and documentation describing terminal output
\usepackage{newunicodechar}
\newunicodechar{✓}{\checkmark}
\newunicodechar{✗}{$\times$}
\newunicodechar{○}{\ensuremath{\circ}}
\newunicodechar{─}{\textemdash}
\newunicodechar{📝}{\texttt{[note]}}
\newunicodechar{🤖}{\texttt{[AI]}}
% LaTeX Error: Unicode character € (U+0080) not set up for use with LaTeX.
\newunicodechar{€}{\texttt{[0x80]}}
% LaTeX Error: Unicode character “ (U+0093) not set up for use with LaTeX.
\newunicodechar{“}{\texttt{[0x93]}}
 after
% \documentclass but before \begin{document}.

% Font setup.  Under pdfLaTeX these are REQUIRED: without [T1]{fontenc} a
% monospace [[code]] quote (or \texttt) cannot select a T1 typewriter shape and
% silently falls back to the proportional OT1 roman font -- i.e. monospace
% "disappears".  iftex guards them so this preamble stays safe when copied into
% a LuaLaTeX/XeLaTeX project (which already use Unicode fonts and need neither).
% Do NOT comment these out.  See latex-writing's references/unicode-and-fonts.md
% ("Symptom 2").
\usepackage{iftex}
\ifpdftex
  \usepackage[utf8]{inputenc}
  \usepackage[T1]{fontenc}
  \usepackage{lmodern}
\fi
\usepackage[swedish,british]{babel}
\usepackage{booktabs}
% The audit tables of `scholar sessions export -f table` are longtables
% with array column types (doc/methods/tables/).
\usepackage{longtable}
\usepackage{array}

\usepackage{noweb}
\noweboptions{shift,breakcode,longxref,longchunks}
% noweb names a chunk by its page and, only when that page holds two or
% more chunks, a letter (\@subpageref tests \2on<page>).  So whether a
% reference such as "papers 1865" reads "1865b" depends on where the
% pages break, and the references' width moves the breaks: the scholar
% document flipped between two layouts on every pass and never settled
% (#279).  With the letter always given, a reference keeps its length
% when its chunk is alone on the page, and the passes settle.
% The "used in" and "continued in" lists (built from \@subpagepair)
% keep noweb's form on purpose: always giving the letter there too made
% a list of more than 140 entries on page 158 change its length with
% the layout (its ranges such as 342-46 and its entries depend on the
% page breaks), and a build from clean still flipped between two
% layouts after 12 passes.
\makeatletter
\def\@subpageref#1#2#3\\{\nwthepagenum{#1}{#2}}
% A run of three or more pages in a "used in" list ("342--46") reaches
% \@nwhyperpagenum with an empty label: noweb.sty's \@pagesl and
% \@pagenumsl build such a range as {lo--hi}{}, and \@nwhyperpagenum
% links every entry to noweb.<label>, so pdfTeX warned "name{noweb.}
% has been referenced but does not exist" (#311).  A range now stays
% unlinked text; an entry with a label links as before, and the text,
% hence the layout, is the same.
\def\@nwhyperpagenum#1#2{%
  \ifx\relax#2\relax #1\else\nwhyperreference{#2}{#1}\fi}
\makeatother

% Verbose citation style: full reference on first use, short form on
% repeats.  This goes with memoir + didactic: didactic sets
% autocite=footnote and places footnotes in the margin, so \autocite
% yields the full reference as a margin note.  Without memoir+didactic
% verbose footnotes land at the page foot; use alphabetic/authoryear
% there instead.  Cite with \autocite in body text; inside \ltnote or any
% \marginpar use \parencite (a footnote inside a margin float is fatal:
% "Float(s) lost").
\usepackage[
  natbib,
  backend=biber,
  style=verbose,
  citestyle=verbose,
  singletitle=false,
  maxbibnames=99,
]{biblatex}
\addbibresource{bibliography.bib}

\usepackage[inline]{enumitem}
\setlist[enumerate]{label=(\arabic*)}

\usepackage[all]{foreign}
\renewcommand{\foreignfullfont}{}
\renewcommand{\foreignabbrfont}{}

%\usepackage{newclude}
\usepackage{import}

\usepackage[strict]{csquotes}
\usepackage[single]{acro}

\usepackage{subcaption}

\usepackage[noend]{algpseudocode}
\usepackage{xparse}

\let\email\texttt

%\usepackage[outputdir=ltxobj]{minted}
\usepackage{minted}
\setminted{autogobble}

\usepackage{pythontex}
\setpythontexoutputdir{.}
\setpythontexworkingdir{..}

\usepackage{fancyvrb}

\usepackage{amsmath}
\usepackage{amssymb}
\usepackage{mathtools}
\usepackage{amsthm}
\usepackage{thmtools}
\usepackage[unq]{unique}
\DeclareMathOperator{\powerset}{\mathcal{P}}

\usepackage[binary-units]{siunitx}

\usepackage{pgf}

\usepackage[colorlinks]{hyperref}
\usepackage[capitalize]{cleveref}
\crefalias{question}{item}
\crefname{question}{question}{questions}
\Crefname{question}{Question}{Questions}
\creflabelformat{question}{#2\textup{#1}#3}

% adapted from https://tex.stackexchange.com/a/30726/17418
\newcommand\chnote[1]{%
  \begingroup
  \renewcommand\thefootnote{}\footnote{#1}%
  \addtocounter{footnote}{-1}%
  \endgroup
}

\usepackage[marginparmargin=right]{didactic}

% Make noweb's hyperlinked identifiers safe in headings (chapter/section titles
% that contain a [[code]] quote).  With -index, such a quote weaves into a
% \nwlinkedidentq hyperlink; in a moving argument (the .toc, the running header,
% the PDF bookmark) hyperref then fails fatally (under newer TeX Live with
% errors like "Undefined control sequence \hyper@@link" / "\Hy@reserved@a", or a
% stack overflow from noweb's self-redefining \nwhyperreference).  Two fixes,
% both required, let the project keep [[...]] in headings instead of falling
% back to \texttt{...}:
%
% (1) Replace \nwhyperreference with noweb's robust hyperref form so it is
%     written literally into moving arguments instead of expanding (and
%     misfiring its terminating \global\let) there, yet still links in the body.
\DeclareRobustCommand\nwhyperreference[2]{\hyperlink{noweb.#1}{#2}}
%
% (2) Strip noweb's code-quote macros from PDF bookmark strings, so a [[code]]
%     quote in a heading yields clean bookmark text without hyperref warnings.
\pdfstringdefDisableCommands{%
  \def\nwlinkedidentq#1#2{#1}%
  \let\Tt\relax
  \let\nwendquote\relax
}

% Unicode character replacements for symbols not in Latin Modern fonts
% These appear in TUI code and documentation describing terminal output
\usepackage{newunicodechar}
\newunicodechar{✓}{\checkmark}
\newunicodechar{✗}{$\times$}
\newunicodechar{○}{\ensuremath{\circ}}
\newunicodechar{─}{\textemdash}
\newunicodechar{📝}{\texttt{[note]}}
\newunicodechar{🤖}{\texttt{[AI]}}
% LaTeX Error: Unicode character € (U+0080) not set up for use with LaTeX.
\newunicodechar{€}{\texttt{[0x80]}}
% LaTeX Error: Unicode character “ (U+0093) not set up for use with LaTeX.
\newunicodechar{“}{\texttt{[0x93]}}


\title{Scholar CLI: A Tool for Structured Literature Searches}
\author{Daniel Bosk\\Dan-Claude van Thropic\thanks{%
    Wrote this document and the Scholar system.
  } \and
  Ric Glassey\\Claude-Eric Anthropicsson\thanks{%
    Wrote Snowball and Tuxedo.
  }
}
\date{\today}

\begin{document}
\frontmatter
\maketitle

\begin{abstract}
A systematic literature review searches several databases with
different interfaces, merges and deduplicates the results, and records
every screening decision.
Scholar is a command-line tool for that work: it unifies access to
bibliographic databases---among them Semantic Scholar, OpenAlex,
DBLP, Web of Science, IEEE Xplore, Scopus, arXiv and HAL---and keeps the
record of every search and decision.

The tool centers on a terminal-based review interface where researchers
navigate papers, read abstracts, record keep/discard decisions with
justifications, and organize findings with themes and notes.
For large result sets, an LLM-assisted classification shows a language
model the user's decisions as examples and records its suggestions for
pending papers as unreviewed until a human confirms or changes them.
Caching, automatic deduplication, and PDF management round out
the workflow support.

This document is the complete literate source of Scholar, weaving together
the design rationale, implementation, and tests into a unified narrative.
The code is extracted (``tangled'') from this document using noweb.
\end{abstract}
\vfill
\inputminted[linenos=false]{text}{../LICENSE}

\clearpage
\tableofcontents*

\mainmatter

\chapter{Introduction}
\label{ch:introduction}

Scholar is a command-line tool for conducting structured literature searches
across multiple academic databases.
It is designed to support systematic literature reviews by providing a unified
interface for searching, reviewing, and documenting the search process.

At its core, Scholar embodies a simple insight: systematic literature reviews
are fundamentally about \emph{traceable decision-making}.
Every paper must be found, evaluated, and either kept or discarded---and each
decision must be justified and reproducible.
Scholar automates the mechanical aspects of this process (searching,
deduplicating, exporting) while preserving human judgment for what matters
most: evaluating relevance.

The tool follows the Unix philosophy of small, composable components.
The command-line interface supports scripting and integration with other tools.
The interactive terminal interface provides efficient navigation for
hands-on review.
And the LLM-assisted classification system extends human judgment without
replacing it---learning from your decisions to suggest (not dictate)
classifications for remaining papers.


\section{Quickstart}

\paragraph{Install}
The recommended installation method is \mintinline{text}{pipx}, which keeps CLI
tools in isolated environments:
\begin{verbatim}
pipx install scholarcli
\end{verbatim}

\paragraph{Configure LLM access (optional)}
The \mintinline{text}{rq} command and LLM-assisted review use the configuration
of Simon Willison's \texttt{llm} package.
To configure providers and select a default model, install the \texttt{llm}
command (or install Scholar with \mintinline{text}{pipx --include-deps} so the
dependency CLI is exposed):
\begin{verbatim}
pipx install llm
# Or: pipx install --include-deps scholarcli
\end{verbatim}

Then configure at least one provider (examples):
\begin{verbatim}
llm install llm-openai-plugin
llm keys set openai

# Or:
llm install llm-anthropic
llm keys set anthropic
\end{verbatim}

Set a default model (used by Scholar unless overridden with \mintinline{text}{--model}):
\begin{verbatim}
llm models
llm models default gpt-4o-mini
\end{verbatim}

\paragraph{Run}
\begin{verbatim}
# Search directly
scholar search "machine learning privacy"

# Start from a research question
scholar rq "How do LLMs support novice programming?" --count 20
\end{verbatim}

\paragraph{Learn by doing}
The interactive tutorials walk through these workflows one command at a
time, checking each step before moving on; the lessons are the chapters of
the Tutorials part of this document.
\begin{verbatim}
scholar tutorial list
scholar tutorial run first-search
\end{verbatim}


\section{Motivation}

Systematic literature reviews require searching multiple databases to ensure 
comprehensive coverage of relevant research.
Each database has its own interface, query syntax, and export format.
Researchers must manually:
\begin{enumerate}
\item Run the same query in multiple databases
\item Export results in various formats
\item Merge and deduplicate results
\item Document which databases were searched, when, and with what queries
\item Review each paper and record decisions with justifications
\end{enumerate}

This manual process is tedious, error-prone, and difficult to reproduce.
Scholar automates these tasks while maintaining the transparency required for 
rigorous systematic reviews.


\section{Features Overview}

Scholar provides the following capabilities:

\subsection{Multi-Database Search}

Search across six academic databases with a single command:
\begin{description}
\item[Semantic Scholar] AI-powered research database covering 200+ million 
  papers across all fields.
  Provides abstracts and open-access PDF links where available.
  Optional API key for higher rate limits.
\item[OpenAlex] Open catalog of 250+ million scholarly works, fully open 
  access.
  Provides comprehensive metadata including abstracts and citation counts.
  Optional email for faster responses.
\item[DBLP] Computer science bibliography with comprehensive coverage of 
  journals and conference proceedings.
  No abstracts, but reliable bibliographic data.
  No authentication required.
\item[Web of Science] Comprehensive citation index covering sciences, social 
  sciences, arts, and humanities.
  Requires API key from Clarivate.
\item[IEEE Xplore] Technical literature from IEEE and partner publishers.
  Strong coverage of electrical engineering, computer science, and related 
  fields.
  Requires API key from IEEE.
\item[arXiv] Preprint repository with strong coverage of physics, mathematics,
  computer science, and related fields.
  No authentication required.
\end{description}

\subsection{Result Management}

\begin{description}
\item[Deduplication] Papers appearing in multiple databases are automatically 
  merged, combining metadata from all sources.
  For example, a paper found in both DBLP (without abstract) and Semantic 
  Scholar (with abstract) appears once with the abstract included.
\item[Caching] Search results are cached locally to avoid redundant API calls 
  and respect rate limits.
  Caches persist across sessions.
\item[Multiple output formats] Results can be displayed as tables (for 
  interactive use), CSV (for spreadsheets), JSON (for scripting), or BibTeX 
  (for \LaTeX{} documents).
\end{description}

\subsection{Interactive Review}

The terminal-based user interface (TUI) supports systematic review workflows:
\begin{description}
\item[Navigation] Browse papers with vim-style keybindings (\texttt{j}/\texttt{k} 
  for up/down).
\item[Viewing] Read full abstracts and metadata for each paper.
\item[Decisions] Mark papers as kept or discarded.
  Discards require motivations (creating an audit trail), while keeps can 
  optionally be tagged with themes for organization.
\item[Notes] Add notes to any paper using your preferred text editor.
  Notes persist across sessions and are included in generated reports.
\item[PDF viewing] Download and view PDFs directly from the TUI.
  PDFs are cached locally for offline access.
\item[Enrichment] Fetch missing abstracts from Semantic Scholar or OpenAlex 
  for papers that lack them.
\item[Sorting and filtering] Order papers by title, year, author, or provider.
  Filter to show only kept, discarded, or pending papers.
\end{description}

\subsection{LLM-Assisted Review}

For large result sets, Scholar can use large language models to assist with 
paper classification:
\begin{description}
\item[Learning from Examples] The LLM learns from papers you have already 
  reviewed and tagged, using them as examples to classify pending papers.
\item[Round-Based Workflow] LLM decisions are marked for human review.
  You can accept, correct, or reject each classification, and corrections 
  become training examples for future rounds.
\item[Confidence Scores] Each LLM decision includes a confidence score 
  (0.0--1.0), allowing you to prioritize reviewing uncertain classifications.
\item[Source Tracking] Decisions track whether they were made by a human, 
  by the LLM and unreviewed, or by the LLM and reviewed.
  This enables progress monitoring and ensures transparency.
\item[Automatic Enrichment] Papers without abstracts are automatically 
  enriched before classification, ensuring the LLM has sufficient context.
\item[Model Flexibility] Uses Simon Willison's \texttt{llm} package, 
  supporting OpenAI, Anthropic, local models, and other providers.
\end{description}

\subsection{Session Management}

Review sessions can be saved and resumed:
\begin{description}
\item[Persistence] Decisions are saved automatically when exiting the TUI.
  Resume a session to continue where you left off.
\item[Named sessions] Use the same session name across multiple searches to 
  accumulate papers for a single review.
\item[Reports] Generate \LaTeX{} reports with BibTeX bibliographies documenting 
  the review process, or export to CSV for spreadsheet analysis.
\end{description}


\section{Typical Workflow}

A typical systematic review workflow with Scholar proceeds as follows:

\paragraph{1. Initial search}
Run a search query and enter the interactive review mode:
\begin{verbatim}
scholar search "privacy-preserving machine learning" \
    --review --name "privacy-ml-review"
\end{verbatim}
The \texttt{--name} flag creates a named session for later resumption.

If you start from a research question rather than a fixed search string, use
\texttt{rq} to generate provider-specific queries with an LLM and save a review
session:
\begin{verbatim}
scholar rq "How can privacy-preserving machine learning be evaluated?" \
    --provider openalex --provider dblp \
    --count 20 --name "privacy-ml-review"
\end{verbatim}

\paragraph{2. Review papers}
In the TUI, browse through papers:
\begin{itemize}
\item Press \texttt{Enter} to view a paper's abstract
\item Press \texttt{K} to keep a paper, or \texttt{T} to keep with themes
\item Press \texttt{d} to discard (you must provide a motivation)
\item Press \texttt{n} to add notes
\item Press \texttt{p} to view the PDF
\item Press \texttt{L} to invoke LLM-assisted classification
\item Press \texttt{q} to quit and save progress
\end{itemize}

\paragraph{3. Add more searches}
Run additional queries and append to the same session:
\begin{verbatim}
scholar search "federated learning privacy" \
    --review --name "privacy-ml-review"
\end{verbatim}
Papers already in the session retain their decisions; new papers are added.

\paragraph{4. Resume later}
Return to an incomplete review:
\begin{verbatim}
scholar sessions resume "privacy-ml-review"
\end{verbatim}

\paragraph{5. Generate reports}
Export the review for inclusion in a paper or thesis:
\begin{verbatim}
scholar sessions export "privacy-ml-review" -f all
\end{verbatim}
This generates \texttt{.tex}, \texttt{.bib}, and \texttt{.csv} files 
documenting all decisions with motivations and notes.
The \LaTeX{} report is a compilable document; with
\texttt{--no-standalone} it becomes a fragment for \verb|\input| in a
paper or thesis (the \texttt{.bib} is written next to it), and
\texttt{-f table} writes an audit table for a backing appendix, one row
per paper with the reason it was kept or excluded:
\begin{verbatim}
scholar sessions export "privacy-ml-review" -f latex --no-standalone -o review
scholar sessions export "privacy-ml-review" -f table --label review -o review-full
\end{verbatim}
The same \texttt{--standalone}/\texttt{--no-standalone} choice applies to
\texttt{scholar llm synthesize --format latex}, whose document ends with
the decision record as an appendix.


\section{Command Reference}

Scholar provides the following commands:

\begin{description}
\item[\texttt{scholar search}] Search bibliographic databases.
  Options control which providers to query, output format, and whether to
  launch the interactive review.
\item[\texttt{scholar rq}] Research-question driven search.
  Uses an LLM to create provider-specific search queries from a research
  question, runs searches with the same filters as \texttt{search}, and
  optionally labels an initial batch of papers.
\item[\texttt{scholar providers}] List available search providers and their
  configuration status (API keys, etc.).
\item[\texttt{scholar sessions}] Manage saved review sessions: list, show, 
  export, or resume.
\item[\texttt{scholar notes}] Browse and manage paper notes: list, show, 
  export, import, or clear.
\item[\texttt{scholar enrich}] Fetch missing abstracts for papers in a saved 
  session.
\item[\texttt{scholar cache}] Manage the search result cache: view statistics, 
  clear, or show the cache directory path.
\item[\texttt{scholar pdf}] Manage PDF downloads: open a PDF by URL, view 
  cache statistics, clear the cache, or show the cache directory.
\end{description}

Each command supports \texttt{--help} for detailed usage information.


\section{Configuration}

Scholar uses environment variables for configuration:

\begin{description}
\item[\texttt{S2\_API\_KEY}] Semantic Scholar API key (optional, increases 
  rate limits).
\item[\texttt{OPENALEX\_EMAIL}] Email for OpenAlex polite pool (optional, 
  faster responses).
\item[\texttt{WOS\_EXPANDED\_API\_KEY} or \texttt{WOS\_STARTER\_API\_KEY}] Web of Science API key (required for WoS
  searches). The legacy \texttt{WOS\_API\_KEY} variable is also supported.
\item[\texttt{IEEE\_API\_KEY}] IEEE Xplore API key (required for IEEE 
  searches).
\item[\texttt{SCHOLAR\_CACHE\_DIR}] Override the default cache directory 
  location.
\item[\texttt{VISUAL} or \texttt{EDITOR}] Text editor for note editing 
  (standard Unix convention).
\end{description}


\section{Document Organization}

This document is a literate program combining documentation and implementation.
It is organized as follows:

\begin{description}
\item[Part I: Core Functionality] The foundation of Scholar: data models for
  papers and search results (\cref{scholar-module}), the caching system that
  avoids redundant API calls (\cref{cache-module}), provider implementations
  for each database (\cref{providers-module}), abstract enrichment from
  secondary sources (\cref{enrich-module}), and PDF downloading and viewing
  (\cref{pdf-module}).
\item[Part II: Review System] The systematic review workflow: persistent notes
  attached to papers (\cref{notes-module}), the review session model that
  tracks decisions and exports reports (\cref{review-module}), and the
  LLM-assisted classification system (\cref{llm-review-module}).
\item[Part III: User Interface] How users interact with Scholar: the
  command-line interface for scripting and quick operations
  (\cref{cli-module}), and the terminal-based user interface for interactive
  paper review (\cref{tui-module}).
\end{description}

Each chapter presents both the design rationale and the implementation, with
tests appearing alongside the code they verify.
The source files are generated (\enquote{tangled}) from this document using
noweb, and can be regenerated at any time from this single source of truth.
A code chunk is named by its page and a letter, and a chunk's name and
every reference to it in the text carry the letter (1865a) even when
the chunk is alone on its page: noweb by itself gives the letter only
when a page holds two or more chunks, which made the page breaks
depend on the references and the references on the page breaks, so
this document never settled.
The lists of where a chunk is used or continued keep noweb's own form,
with the letter only on such pages.


\part{Tutorials}

The chapters in this part are the literate sources of the interactive
tutorials that \texttt{scholar tutorial} runs.
Each tangles to the Markdown that pytorial executes and weaves to a
chapter that explains how the lesson teaches: which aspect of Scholar a
step lets the learner vary, what it holds invariant, and why the steps
come in this order.
That design reasoning is the object of these chapters, so it stands in
the body text; the learner never sees it, because only the Markdown is
shipped.
The tutorials come first because they are the shortest path from a fresh
install to a finished review; the parts that follow explain how the tool
does what these chapters teach.

\input{../src/scholar/tutorials/first-search.tex}
\input{../src/scholar/tutorials/reviewing-in-the-tui.tex}
\input{../src/scholar/tutorials/enrich-verify-export.tex}
\input{../src/scholar/tutorials/tiered-search.tex}
\input{../src/scholar/tutorials/llm-synthesis.tex}
\input{../src/scholar/tutorials/scholar-with-an-agent.tex}

\part{Core Functionality}

\input{../src/scholar/scholar.tex}
\input{../src/scholar/utils.tex}
\input{../src/scholar/cache.tex}
\input{../src/scholar/ratelimit.tex}
\input{../src/scholar/network.tex}

\input{../src/scholar/providers.tex}
\input{../src/scholar/enrich.tex}
\input{../src/scholar/crossref.tex}
\input{../src/scholar/unpaywall.tex}
\input{../src/scholar/europepmc.tex}
\input{../src/scholar/preprint.tex}
\input{../src/scholar/pdf.tex}

\part{Review System}

\input{../src/scholar/notes.tex}
\input{../src/scholar/review.tex}
\input{../src/scholar/llm_review.tex}
\input{../src/scholar/agreement.tex}

\part{User Interface}

\input{../src/scholar/cli.tex}
\input{../src/scholar/tui.tex}

\backmatter

\end{document}
